\section{Introducción}
\label{sec:introduccion}

Acortar una grabación para que entre en una duración fija obliga a decidir qué parte se pierde. Truncar corta el final de la pieza y descarta justamente el cierre. Acelerar la reproducción de manera uniforme conserva todo el material pero altera el tempo y la altura de principio a fin. El recorte consciente del contenido \cite{wenner2013} evita ambos problemas descartando tramos internos y uniendo sus extremos, con la apuesta de que la música se repite y de que un salto entre dos puntos parecidos pasa inadvertido. Esa es la técnica con la que se producen los \emph{radio edits} (las versiones acortadas de un tema para emisión radial) y la que incorporan los editores de audio.

El problema que resolvemos en este trabajo es el de elegir dónde cortar. Dada una grabación representada como una secuencia de $n$ pulsos, cada uno descripto por un vector de características $f_i \in \mathbb{R}^d$, se pide conservar exactamente $k$ de ellos, siempre el primero y el último, de modo que la suma de los costos de los saltos entre pulsos conservados consecutivos sea mínima, donde el costo de un salto es la distancia euclídea entre el pulso que debería haber sonado y el que efectivamente suena.

Implementamos en C++ tres algoritmos exactos y comparamos hasta dónde llega cada uno. Fuerza bruta enumera todas las selecciones posibles y se queda con la más barata. Backtracking recorre el mismo espacio pero abandona cada rama apenas una poda por factibilidad o una por optimalidad muestra que no puede mejorar lo ya encontrado. Programación dinámica resuelve el problema por subproblemas sobre una tabla y reconstruye la selección óptima a partir de ella. Reimplementamos backtracking y programación dinámica en Python con asistencia de una herramienta de inteligencia artificial para comparar lenguajes con el mismo algoritmo. La experimentación usa tres familias de instancias sintéticas diseñadas para separar el efecto de la estructura de los datos, más grabaciones reales (tres temas instrumentales, tres escenas de diálogo y un tramo de radioteatro, de 183 a 645 pulsos) cuyos recortes reconstruimos en audio.

Los tres algoritmos devuelven el mismo costo óptimo en las 93 instancias sintéticas y en los 27 prefijos de un tema real, pero solo dos de ellos escalan al tamaño de una grabación. Con $k = n/2$, fuerza bruta tarda 21,3 s en $n = 32$ y cada pulso adicional duplica ese tiempo, mientras que backtracking y programación dinámica resuelven la misma instancia en menos de 0,04 ms. Sobre un tema real de 497 pulsos la programación dinámica en C++ tarda 43,8 ms, mientras que backtracking depende del contenido y llega a ser 45 veces más lento en un prefijo intermedio. La misma programación dinámica en Python tarda entre 104 y 110 veces más que en C++. En audio real la selección óptima cuesta en la mediana un 30\,\% de lo que cuesta truncar y concentra el descarte en uno o dos tramos largos. La recomendación del grupo para la práctica es la programación dinámica en C++, con la versión en Python como referencia cruzada de corrección.
